\documentclass[9pt,aspectratio=169,xcolor=table]{beamer}

%~~~~~~~~~~~~~~~~~~~~~~~~~~~~~~~~~~~~~~~~~~~~~~~~~~~~~~~~~~~~~~~~~~~~~~~~~~~~~~
% Use roboto Font (recommended)
\usepackage[sfdefault]{roboto}
\usepackage[utf8]{inputenc}
\usepackage[T1]{fontenc}
%~~~~~~~~~~~~~~~~~~~~~~~~~~~~~~~~~~~~~~~~~~~~~~~~~~~~~~~~~~~~~~~~~~~~~~~~~~~~~~

%~~~~~~~~~~~~~~~~~~~~~~~~~~~~~~~~~~~~~~~~~~~~~~~~~~~~~~~~~~~~~~~~~~~~~~~~~~~~~~
% Define where theme files are located. ('/styles')
\usepackage{styles/fluxmacros}
\usefolder{styles}
% Use Flux theme v0.1 beta
% Available style: asphalt, blue, red, green, gray
% "red" is the closest built-in style to maroon; its exact colours are
% overridden with a true maroon palette just below, after the theme loads.
\usetheme[style=red]{flux}
%~~~~~~~~~~~~~~~~~~~~~~~~~~~~~~~~~~~~~~~~~~~~~~~~~~~~~~~~~~~~~~~~~~~~~~~~~~~~~~

%~~~~~~~~~~~~~~~~~~~~~~~~~~~~~~~~~~~~~~~~~~~~~~~~~~~~~~~~~~~~~~~~~~~~~~~~~~~~~~
% SKEE1033 maroon palette override.
% The flux theme's colortheme defines primaryLight/primary/secondary/tertiary
% as plain \definecolor calls (not \providecolor), so simply redefining them
% here, AFTER \usetheme, wins -- xcolor allows redefinition and every
% \setbeamercolor in beamercolorthemeflux.sty was already evaluated against
% these names, so the redefinition propagates everywhere automatically.
\definecolor{primaryLight}{HTML}{9B2242} % lighter maroon -- headers, accents
\definecolor{primary}{HTML}{7B1E3A}      % core maroon -- titles, block headers
\definecolor{secondary}{HTML}{3A5A6B}    % muted teal -- complementary accent
\definecolor{tertiary}{HTML}{8C6A3F}     % warm bronze -- example/third accent
%~~~~~~~~~~~~~~~~~~~~~~~~~~~~~~~~~~~~~~~~~~~~~~~~~~~~~~~~~~~~~~~~~~~~~~~~~~~~~~

%~~~~~~~~~~~~~~~~~~~~~~~~~~~~~~~~~~~~~~~~~~~~~~~~~~~~~~~~~~~~~~~~~~~~~~~~~~~~~~
% Extra packages
\usepackage{booktabs}
\usepackage{colortbl}
\usepackage{ragged2e}
\usepackage{amssymb}
\usepackage{multirow}
\usepackage{tikz}
\usetikzlibrary{arrows.meta, positioning, fit, backgrounds, calc, shapes.geometric}
\setbeamertemplate{itemize items}[square]
\usepackage[scaled=.92]{beramono}
\usepackage{listings}
\usepackage{array}
% Table striping: odd rows white, even rows very light gray.
\definecolor{stripe}{gray}{0.94}
\newcommand{\striped}{\rowcolors{1}{white}{stripe}}
%~~~~~~~~~~~~~~~~~~~~~~~~~~~~~~~~~~~~~~~~~~~~~~~~~~~~~~~~~~~~~~~~~~~~~~~~~~~~~~

%~~~~~~~~~~~~~~~~~~~~~~~~~~~~~~~~~~~~~~~~~~~~~~~~~~~~~~~~~~~~~~~~~~~~~~~~~~~~~~
% Code listing style (lightweight analogue of the book's skpython style).
% This course is Python-only, so Python is set as the lstlisting default via
% \lstset below -- every \begin{lstlisting} needs no language=... of its own.
\definecolor{codebg}{HTML}{FBF2F4}
\definecolor{codegreen}{rgb}{0,0.45,0}
\definecolor{codestring}{rgb}{0.58,0.1,0.2}
\lstdefinestyle{skpython}{
    basicstyle=\ttfamily\footnotesize,
    keywordstyle=\color{primary}\bfseries,
    commentstyle=\itshape\color{codegreen},
    stringstyle=\ttfamily\color{codestring},
    backgroundcolor=\color{codebg},
    breaklines=true,
    keepspaces=true,
    showspaces=false,
    showstringspaces=false,
    frame=single,
    rulecolor=\color{primary!40},
    numbers=none,
}
\lstset{language=Python, style=skpython}
%~~~~~~~~~~~~~~~~~~~~~~~~~~~~~~~~~~~~~~~~~~~~~~~~~~~~~~~~~~~~~~~~~~~~~~~~~~~~~~

%~~~~~~~~~~~~~~~~~~~~~~~~~~~~~~~~~~~~~~~~~~~~~~~~~~~~~~~~~~~~~~~~~~~~~~~~~~~~~~
% TikZ block styles for the computation-cycle and control-structure diagrams
\definecolor{cycfill}{HTML}{F3DCE4}
\definecolor{cycline}{HTML}{7B1E3A}
\definecolor{cycdofill}{HTML}{E8EEF0}
\definecolor{cycdoline}{HTML}{3A5A6B}
\tikzset{
  cycstep/.style = {font=\sffamily\small, align=center, rounded corners=2pt,
                     draw=cycline, fill=cycfill, line width=0.7pt,
                     minimum height=7mm, minimum width=20mm, inner sep=3pt},
  cycdo/.style   = {cycstep, fill=cycdofill, draw=cycdoline},
  flow/.style    = {-{Stealth[length=2mm]}, thick, draw=cycline},
}
%~~~~~~~~~~~~~~~~~~~~~~~~~~~~~~~~~~~~~~~~~~~~~~~~~~~~~~~~~~~~~~~~~~~~~~~~~~~~~~

%%%%%%%%%%%%%%%%%%%%%%%%%%%%%%%%%%%
%% SECTION SEPARATOR (ToC recap) %%
%%%%%%%%%%%%%%%%%%%%%%%%%%%%%%%%%%%
\AtBeginSection[]
{
{
    \setbeamercolor{background canvas}{bg=primaryLight!6}
    \begin{frame}[plain]
        \frametitle{Table of Contents}
        \tableofcontents[currentsection]
    \end{frame}
}
}
%%%%%%%%%%%%%%%%%%%%%%%%%%%%%%%%%%%

%~~~~~~~~~~~~~~~~~~~~~~~~~~~~~~~~~~~~~~~~~~~~~~~~~~~~~~~~~~~~~~~~~~~~~~~~~~~~~~
% Information
\title{Handling Imperfect Engineering Data}
\subtitle{\emph{Engineering Computation with Python} --- Chapter 11}
\author{Mun\'{i}m Zabidi}
\institute{\color{primary}Faculty of Electrical Engineering, Universiti Teknologi Malaysia}
\date{\today}
\titlegraphic{assets/logoutmwhite.png}
%~~~~~~~~~~~~~~~~~~~~~~~~~~~~~~~~~~~~~~~~~~~~~~~~~~~~~~~~~~~~~~~~~~~~~~~~~~~~~~

\begin{document}

\titlepage

\begin{frame}{Table of Contents}
    \tableofcontents
\end{frame}

%===============================================================================
\section{Errors, Exceptions and try/except}
%===============================================================================

\begin{frame}{Two Different Things Go Wrong}
    \leavevmode
    This chapter deals with both.

    \vspace{2mm}
    \begin{columns}[t]
        \begin{column}{.5\textwidth}
            \begin{block}{An error in your \alert{code}}
                \small A file that does not exist, a division by zero, a
                value of the wrong type. Python detects these
                \alert{immediately} and stops the program with an
                exception --- unless you handle it.
            \end{block}
        \end{column}
        \begin{column}{.5\textwidth}
            \begin{alertblock}{A flaw in your \alert{data}}
                \small A sensor dropout, a power-spike glitch, a
                physically implausible value. Python does \alert{not}
                detect these --- a corrupted reading is still a
                perfectly valid floating-point number.
            \end{alertblock}
        \end{column}
    \end{columns}

    \vspace{3mm}
    \begin{block}{The job split}
        A robust program \alert{survives} the first kind of problem. A
        careful engineer \alert{investigates and repairs} the second.
    \end{block}
\end{frame}

\begin{frame}{What Is an Error?}
    \leavevmode
    In Python, errors fall into two broad categories:

    \vspace{3mm}
    \begin{columns}[t]
        \begin{column}{.5\textwidth}
            \begin{block}{Syntax errors}
                \small Caught by Python \alert{before} the program even
                runs. Code violates the grammatical rules of the
                language.
            \end{block}
        \end{column}
        \begin{column}{.5\textwidth}
            \begin{block}{Exceptions}
                \small Occur \alert{while} the program is running.
                Something unexpected happens during execution, even
                though the code is syntactically correct.
            \end{block}
        \end{column}
    \end{columns}

    \vspace{4mm}
    \begin{alertblock}{Key fact}
        Exceptions are a \alert{normal} part of real programs --- not
        rare bugs. Any program that interacts with users, files, or
        external systems should expect them and be designed to handle
        them.
    \end{alertblock}
\end{frame}

\begin{frame}[fragile]{The try / except Anatomy}
    \leavevmode
    The basic structure for handling exceptions has four parts:

\begin{lstlisting}
try:
    # risky code
except ErrorType:
    # handle the error
else:
    # runs if no error
finally:
    # always runs (cleanup)
\end{lstlisting}

    \vspace{1mm}
    \begin{itemize}\small
        \item \alert{try} --- code that might raise an exception.
        \item \alert{except} --- catches and handles a specific
              exception type, if one occurs.
        \item \alert{else} --- runs only if \texttt{try} completed with
              no exception.
        \item \alert{finally} --- always runs, exception or not; the
              right place for cleanup.
    \end{itemize}
\end{frame}

\begin{frame}{What Is an Exception?}
    \leavevmode
    An exception is a runtime error object created when Python cannot
    continue executing a piece of code normally.

    \vspace{2mm}
    \begin{columns}[t]
        \begin{column}{.5\textwidth}
            \begin{block}{Not handled}
                \small The program \alert{crashes}.
            \end{block}
        \end{column}
        \begin{column}{.5\textwidth}
            \begin{block}{Handled}
                \small The program \alert{recovers and continues}.
            \end{block}
        \end{column}
    \end{columns}

    \vspace{3mm}
    \textbf{Typical situations that raise exceptions:}
    \begin{itemize}
        \item Dividing by zero.
        \item Converting text to a number when the text is not a valid
              number.
        \item Opening a file that does not exist.
    \end{itemize}
\end{frame}

%===============================================================================
\section{Common Exception Types}
%===============================================================================

\begin{frame}{Common Exception Types: The Full Roster}
    \leavevmode
    \small
    Python defines many built-in exception types. Each one tells you
    specifically what went wrong.

    \vspace{2mm}
    \striped
    \begin{center}
    \scriptsize
    \begin{tabular}{@{}ll@{}}
    \toprule
    \textbf{Exception} & \textbf{Meaning} \\
    \midrule
    \texttt{ValueError} & Correct type, invalid or inappropriate value \\
    \texttt{TypeError} & Operation on an incompatible/unsupported type \\
    \texttt{ZeroDivisionError} & Division (or modulo) by zero \\
    \texttt{IndexError} & Sequence index out of range \\
    \texttt{KeyError} & Dictionary key does not exist \\
    \texttt{FileNotFoundError} & File or path could not be found \\
    \texttt{ImportError} & Module or name cannot be found or loaded \\
    \texttt{AttributeError} & Object has no such attribute/method \\
    \texttt{NameError} & Name has not been defined \\
    \texttt{SyntaxError} & Code violates Python's grammar rules \\
    \texttt{IndentationError} & Incorrect or missing indentation \\
    \texttt{RuntimeError} & Generic runtime failure (e.g.\ infinite recursion) \\
    \bottomrule
    \end{tabular}
    \end{center}

    \vspace{2mm}
    {\scriptsize The next few slides look closely at the five you will
    meet most often. The rest appear in the reference table at the end
    of this section.}
\end{frame}

\begin{frame}[fragile]{ValueError}
    \leavevmode
    Correct \alert{type}, but an invalid or inappropriate \alert{value}.

    \begin{columns}[t]
        \begin{column}{.5\textwidth}
\begin{lstlisting}
x = int("abc")
print(x)
\end{lstlisting}
            \texttt{\scriptsize ValueError: invalid literal for int()}
        \end{column}
        \begin{column}{.5\textwidth}
            \small
            \textbf{Common causes:}
            \begin{itemize}\scriptsize
                \item Converting invalid values
                \item Passing an invalid value to a function
                \item Incorrect input format
                \item Values outside the expected range
            \end{itemize}
        \end{column}
    \end{columns}

    \vspace{2mm}
    \begin{block}{How to fix}
        \small Ensure the value matches the expected format; validate
        user input before conversion; use \texttt{try}/\texttt{except}.
    \end{block}
\end{frame}

\begin{frame}[fragile]{TypeError and ZeroDivisionError}
    \leavevmode
    \begin{columns}[t]
        \begin{column}{.5\textwidth}
            \textbf{TypeError} --- incompatible data types.
\begin{lstlisting}
x = 10
y = "5"
print(x + y)
\end{lstlisting}
            \texttt{\scriptsize TypeError: unsupported operand type(s)}
        \end{column}
        \begin{column}{.5\textwidth}
            \textbf{ZeroDivisionError} --- division by zero.
\begin{lstlisting}
x = 10
y = 0
print(x / y)
\end{lstlisting}
            \texttt{\scriptsize ZeroDivisionError: division by zero}
        \end{column}
    \end{columns}

    \vspace{2mm}
    \begin{block}{How to fix}
        \small Convert types properly (\texttt{int()}, \texttt{str()},
        \texttt{float()}); check a value is nonzero before dividing;
        use \texttt{type()} to debug.
    \end{block}
\end{frame}

\begin{frame}[fragile]{IndexError and KeyError}
    \leavevmode
    \begin{columns}[t]
        \begin{column}{.5\textwidth}
            \textbf{IndexError} --- sequence index out of range.
\begin{lstlisting}
nums = [1, 2, 3]
print(nums[5])
\end{lstlisting}
            \texttt{\scriptsize IndexError: list index out of range}
        \end{column}
        \begin{column}{.5\textwidth}
            \textbf{KeyError} --- dictionary key does not exist.
\begin{lstlisting}
data = {"name": "Amanda"}
print(data["email"])
\end{lstlisting}
            \texttt{\scriptsize KeyError: 'email'}
        \end{column}
    \end{columns}

    \vspace{2mm}
    \begin{block}{How to fix}
        \small Check length with \texttt{len()} before indexing; check
        key existence with \texttt{in}, or use \texttt{.get()} to avoid
        the error entirely.
    \end{block}
\end{frame}

\begin{frame}[fragile]{FileNotFoundError}
    \leavevmode
    Python cannot find or access the requested file.

\begin{lstlisting}
file = open("data.txt", "r")
content = file.read()
file.close()
\end{lstlisting}
    \texttt{FileNotFoundError: No such file or directory}

    \vspace{3mm}
    \textbf{Common causes:}
    \begin{itemize}
        \item File does not exist, or wrong name/extension.
        \item Incorrect file path, or file in a different directory.
    \end{itemize}

    \vspace{2mm}
    \begin{block}{How to fix}
        \small Check the file name and extension; verify the path;
        confirm the file exists; use an absolute path if needed.
    \end{block}
\end{frame}

\begin{frame}{The Rest of the Roster (Reference)}
    \leavevmode
    \small
    Six more exception types appear in the book, each with the same
    pattern of common causes and fix. Quick summary:

    \vspace{2mm}
    \begin{itemize}\scriptsize
        \item \alert{ImportError} --- a module or name cannot be found/loaded
              (e.g.\ a typo: \texttt{from math import sqr}).
        \item \alert{AttributeError} --- an object has no such
              attribute/method (e.g.\ \texttt{(10).append(5)}).
        \item \alert{NameError} --- a variable or function was never
              defined, or is misspelled.
        \item \alert{SyntaxError} --- caught before the program even
              runs; a missing colon, bracket, or quote.
        \item \alert{IndentationError} --- incorrect or missing
              spacing; Python uses indentation to define structure.
        \item \alert{RuntimeError} --- a generic failure during
              execution, e.g.\ infinite recursion.
    \end{itemize}

    \vspace{2mm}
    \begin{alertblock}{A related but distinct problem}
        A \textbf{logical error} raises no exception at all --- the
        code runs, but produces the \alert{wrong} output. These are
        often the hardest bugs to find, because Python gives no error
        message.
    \end{alertblock}
\end{frame}

%===============================================================================
\section{Writing Robust Code}
%===============================================================================

\begin{frame}[fragile]{Handling User Input Safely}
    \leavevmode
    A very common source of exceptions is user input:

\begin{lstlisting}
x = int(input("Enter a number"))
\end{lstlisting}

    \vspace{1mm}
    This line can fail in more than one way:
    \begin{itemize}
        \item The user types text instead of a number
              $\rightarrow$ \texttt{ValueError}.
        \item The user presses Enter without typing anything
              $\rightarrow$ \texttt{ValueError}.
    \end{itemize}

    \vspace{2mm}
    \begin{block}{Multiple except clauses}
        Different problems need different handlers --- several
        \texttt{except} clauses can attach to one \texttt{try}, so the
        program responds to each kind of error without crashing.
    \end{block}
\end{frame}

\begin{frame}[fragile]{Handling Multiple Exceptions at Once}
    \leavevmode
    The underlying lesson: programs must not blindly trust their
    users.

\begin{lstlisting}
try:
    x = int(input("Enter a number: "))
    print(10 / x)
except ValueError:
    print("Please enter a valid number")
except ZeroDivisionError:
    print("Cannot divide by zero")
\end{lstlisting}

    \vspace{2mm}
    \begin{block}{Files are external resources too}
        They may not exist, paths may be wrong, permissions may be
        denied. Exception handling around file operations protects
        against corrupted or locked files and prevents data loss.
    \end{block}
\end{frame}

\begin{frame}[fragile]{finally: Cleanup No Matter What}
    \leavevmode
    The file may fail to open, or fail partway through reading. The
    \texttt{finally} block runs regardless --- the right place to close
    a file.

\begin{lstlisting}
try:
    f = open("data.txt", "r")
    content = f.read()
except FileNotFoundError:
    print("File not found!")
except IOError:
    print("Error reading file!")
finally:
    print("Closing file...")
    try:
        f.close()
    except:
        pass
\end{lstlisting}

    {\scriptsize The nested \texttt{try} inside \texttt{finally} guards
    against the case where \texttt{f} was never opened in the first
    place.}
\end{frame}

%===============================================================================
\section{Floating-Point Pitfalls}
%===============================================================================

\begin{frame}[fragile]{The Problem with Floating-Point Comparisons}
    \leavevmode
    Some errors are silent --- no exception is raised, yet the result
    is wrong.

\begin{lstlisting}
a = 0.1 + 0.2
if a == 0.3:
    print("a is 0.3")
\end{lstlisting}

    \vspace{1mm}
    \begin{alertblock}{This print statement never executes}
        \texttt{0.1 + 0.2} does not exactly equal \texttt{0.3} in
        floating-point arithmetic.
    \end{alertblock}

    \vspace{2mm}
    \textbf{Why:} floats are stored in binary, not decimal; many
    decimal numbers cannot be represented exactly; \texttt{==} checks
    for an \alert{exact} match.
\end{frame}

\begin{frame}[fragile]{Epsilon-Based Comparison}
    \leavevmode
    Compare with a \alert{tolerance} rather than exact equality:
    \[
    |a - b| < \varepsilon \implies \text{treat } a \text{ and } b
    \text{ as equal}
    \]

\begin{lstlisting}
epsilon = 1e-9  # acceptable error margin
a = 0.1 + 0.2
b = 0.3
if abs(a - b) < epsilon:
    print("Equal within epsilon")
\end{lstlisting}

    \vspace{2mm}
    {\small This works, but it is easy to misuse --- e.g.\ choosing an
    epsilon inappropriate for the magnitude of the numbers compared.}
\end{frame}

\begin{frame}[fragile]{The Recommended Method: math.isclose()}
    \leavevmode
    Python's standard library handles this automatically --- safer and
    clearer than writing manual checks.

\begin{lstlisting}
import math
print(math.isclose(0.1 + 0.2, 0.3))              # True
print(math.isclose(1.0001, 1.0))                 # False
print(math.isclose(1.0001, 1.0, rel_tol=1e-3))   # True
\end{lstlisting}

    \vspace{2mm}
    \begin{itemize}
        \item \texttt{rel\_tol} --- relative tolerance (a fraction of
              the larger value).
        \item \texttt{abs\_tol} --- absolute tolerance, useful near
              zero.
    \end{itemize}
\end{frame}

\begin{frame}{Error Handling: Key Takeaways}
    \leavevmode
    \small
    \begin{itemize}
        \item Errors are a normal part of real programs.
        \item \texttt{try}/\texttt{except}/\texttt{finally} prevents
              crashes and protects resources such as open files.
        \item Catching a \alert{specific} exception type is far better
              practice than a bare \texttt{except:}, which hides
              faults you did not anticipate.
        \item Floating-point numbers are inexact --- avoid \texttt{==};
              use \texttt{math.isclose()}.
    \end{itemize}

    \vspace{3mm}
    \begin{center}
        \large\itshape\color{primary}
        ``Good programs don't avoid errors --- they handle them.''
    \end{center}
\end{frame}

%===============================================================================
\section{Missing Data}
%===============================================================================

\begin{frame}[fragile]{The Anatomy of Missing Data: NaN}
    \leavevmode
    A missing value in a NumPy array of floats is represented by
    \alert{NaN} (Not-a-Number), an IEEE 754 special value.

\begin{lstlisting}
import numpy as np

x = np.array([1.0, 2.0, np.nan, 4.0])

print(x + 1)          # [2. 3. nan 5.]  -- NaN spreads
print(x == x)          # [True True False True] -- NaN != NaN
print(np.mean(x))      # nan -- one bad value poisons the result
\end{lstlisting}

    \vspace{1mm}
    \begin{alertblock}{The practical danger}
        A single dropped sensor reading, left unhandled, silently
        turns an entire statistical summary into \texttt{nan}.
    \end{alertblock}
\end{frame}

\begin{frame}{Should This Gap Even Be Filled?}
    \leavevmode
    \small
    Before reaching for \texttt{np.interp()}, pause: the right answer
    depends on what the gap \alert{represents}.

    \vspace{2mm}
    \begin{itemize}
        \item \textbf{Short gap, smooth signal?} Usually safe to
              \alert{interpolate}.
        \item \textbf{Isolated entry, not a time series?} No
              "neighbours in time" --- a \alert{mean, median, or
              model-based estimate} may fit better.
        \item \textbf{Large fraction missing?} Filling it manufactures
              more signal than was measured --- \alert{discard the
              segment} or investigate the sensor instead.
        \item \textbf{Is "missing" itself meaningful?} Then
              \alert{keep it as missing and document why}.
    \end{itemize}

    \vspace{2mm}
    \begin{block}{Notice}
        Three of these four branches do \alert{not} end in
        interpolation. Which tool to reach for matters more than
        knowing any one tool well.
    \end{block}
\end{frame}

\begin{frame}[fragile]{Logical Masking and Filtering}
    \leavevmode
    Before a gap can be repaired, it must first be \alert{located}.
    \texttt{np.isnan()} returns a Boolean array, \texttt{True} wherever
    a value is \texttt{NaN}:

\begin{lstlisting}
raw_data = np.array([12.5, np.nan, 13.1, 12.9, np.nan, 13.4])

is_missing = np.isnan(raw_data)
print(is_missing)  # [False True False False True False]

clean_data = raw_data[~np.isnan(raw_data)]
print(clean_data)  # [12.5 13.1 12.9 13.4]
\end{lstlisting}

    \vspace{2mm}
    {\small The \texttt{\textasciitilde} operator flips each Boolean,
    turning "is missing" into "is not missing". Filtering shortens the
    array and discards the \alert{position} of the missing readings ---
    for a time series, that position is often exactly what you need.}
\end{frame}

\begin{frame}[fragile]{Repairing a Signal by Interpolation}
    \leavevmode
    \texttt{np.interp()} estimates each missing value from its
    neighbours, preserving the full length and timing of the signal:

\begin{lstlisting}
np.interp(x[nans], x[~nans], y[~nans])
\end{lstlisting}

    \vspace{1mm}
    \begin{itemize}\small
        \item \texttt{x[nans]} --- positions where a value is missing.
        \item \texttt{x[\textasciitilde nans]} --- positions with a
              good value.
        \item \texttt{y[\textasciitilde nans]} --- the good values
              themselves.
    \end{itemize}

    \vspace{2mm}
    \begin{alertblock}{EE Focus: Patching a Sensor Dropout}
        A temperature sensor logging twice a second suffers a
        3-second dropout (loose connector). Seven consecutive readings
        are lost.
    \end{alertblock}
\end{frame}

\begin{frame}[fragile]{Worked Example: Sensor Dropout Repair}
    \leavevmode
\begin{lstlisting}
t = np.arange(0, 20, 0.5)
temp_true = 25 + 3 * np.sin(0.3 * t)

temp_measured = temp_true.copy()
dropout = (t >= 8) & (t <= 11)
temp_measured[dropout] = np.nan

nans = np.isnan(temp_measured)
temp_repaired = temp_measured.copy()
temp_repaired[nans] = np.interp(t[nans], t[~nans],
                                 temp_measured[~nans])
\end{lstlisting}
    \texttt{\scriptsize Number of missing samples: 7 \quad Max repair error: 0.165 deg C}

    \vspace{2mm}
    \begin{block}{But not a universal rule}
        Interpolation assumes the signal changes \alert{smoothly}
        across the gap. If a step change happened inside the missing
        interval, a straight line can manufacture a plausible but
        false trajectory.
    \end{block}
\end{frame}

%===============================================================================
\section{Outliers}
%===============================================================================

\begin{frame}{What Is an Outlier?}
    \leavevmode
    A missing value is easy to define: the sensor recorded nothing. An
    outlier is trickier --- the sensor \alert{did} record something;
    the question is whether to trust it.

    \vspace{2mm}
    \begin{alertblock}{A statistical description, not a verdict}
        An outlier is a measurement unusually far from the surrounding
        data. That alone does not tell you whether it is a fault or
        the most important sample in the dataset.
    \end{alertblock}

    \vspace{2mm}
    {\small Distinguishing fault from genuine event is far more
    important --- and far harder --- than choosing a detection
    formula.}
\end{frame}

\begin{frame}{Three Kinds of Suspicious Reading}
    \leavevmode
    \small
    \begin{block}{Type A --- Obvious measurement error}
        A temperature trace near $26.5\,^{\circ}\text{C}$ shows one
        sample at $91.3\,^{\circ}\text{C}$, then returns to
        $26.7\,^{\circ}\text{C}$. Nothing heats and cools that fast ---
        almost certainly a glitch.
    \end{block}

    \vspace{1mm}
    \begin{block}{Type B --- Suspicious but uncertain}
        One sample at $31.3\,^{\circ}\text{C}$. Could be a small
        glitch, or the leading edge of a real transient. \alert{Do not
        correct automatically.}
    \end{block}

    \vspace{1mm}
    \begin{block}{Type C --- Genuine event}
        $26 \to 27 \to 28 \to 45 \to 43 \to 35 \to 29\,^{\circ}\text{C}$
        --- rises and falls smoothly. Very likely the most important
        measurement in the dataset.
    \end{block}
\end{frame}

\begin{frame}{Never Automatically Delete or Replace}
    \leavevmode
    A statistical test (Z-score, IQR --- Chapter~\ref{ch:data-analytics})
    flags Type A, B, and C readings \alert{equally well}, because
    statistically they all look the same: unusually far from the rest
    of the data.

    \vspace{3mm}
    \begin{alertblock}{Only engineering judgement tells them apart}
        Looking at neighbouring samples, the physical system, and
        other measured channels is what distinguishes a fault from a
        real event --- a formula alone cannot.
    \end{alertblock}

    \vspace{3mm}
    \begin{center}
        \Large\color{primary}\textbf{Detect $\to$ Flag $\to$
        Investigate $\to$ Decide $\to$ Document}
    \end{center}
\end{frame}

\begin{frame}[fragile]{Worked Example: Investigating a Suspicious Reading}
    \leavevmode
    The bench temperature log shows one sample far above its
    neighbours:

\begin{lstlisting}
t = np.arange(0, 30, 1.0)
temp = 25 + 1.5*np.sin(0.2*t) + np.random.normal(0, 0.1, t.size)
temp[18] = 45.8  # a single sample looks far too high
\end{lstlisting}

    \vspace{1mm}
    \textbf{The investigation, step by step:}
    \begin{enumerate}\small
        \item \alert{Plot it} --- the point clearly stands apart.
        \item \alert{Compare with neighbours} --- no gradual rise or
              fall leading up to it.
        \item \alert{Check physical plausibility} --- is a
              $20\,^{\circ}\text{C}$ jump in one second consistent
              with the board's thermal mass? Usually not.
        \item \alert{Check other channels} --- are voltage and current
              normal at sample 18?
    \end{enumerate}
\end{frame}

\begin{frame}[fragile]{Only Then: Flag and Repair}
    \leavevmode
    Only once the reasoning points to a fault does it make sense to
    flag and repair --- using the same tool already used for missing
    data:

\begin{lstlisting}
is_suspect = np.zeros(t.size, dtype=bool)
is_suspect[18] = True

temp_flagged = temp.copy()
temp_flagged[is_suspect] = np.nan

temp_repaired = temp_flagged.copy()
temp_repaired[is_suspect] = np.interp(
    t[is_suspect], t[~is_suspect], temp_flagged[~is_suspect])
\end{lstlisting}

    \vspace{2mm}
    \begin{alertblock}{The conclusion could reverse}
        If voltage and current also rose at sample 18, this is a
        genuine Type C event --- interpolating it away would erase the
        most important measurement in the log.
    \end{alertblock}
\end{frame}

\begin{frame}{Data Imputation: Three Strategies}
    \leavevmode
    \small
    \textbf{Imputation} --- replacing a missing value with an estimate.
    Interpolation is one way; it is not the only way.

    \vspace{2mm}
    \begin{columns}[t]
        \begin{column}{.33\textwidth}
            \begin{block}{Mean / median}
                \scriptsize Fast and simple. Good when missing values
                are scattered, not contiguous. A median resists
                outliers better than a mean.
            \end{block}
        \end{column}
        \begin{column}{.33\textwidth}
            \begin{block}{Forward-fill}
                \scriptsize Repeats the last known good value.
                Suits signals that change only occasionally (a status
                flag, a calibration value).
            \end{block}
        \end{column}
        \begin{column}{.33\textwidth}
            \begin{block}{Interpolation}
                \scriptsize Estimates from neighbours. Suits a
                smoothly varying signal with a short gap.
            \end{block}
        \end{column}
    \end{columns}

    \vspace{3mm}
    {\small There is no universally "correct" choice --- it depends on
    how the underlying quantity actually behaves between
    measurements.}
\end{frame}

\begin{frame}{Quick Reference: Handling Imperfect Data}
    \leavevmode
    \small
    \striped
    \begin{center}
    \scriptsize
    \begin{tabular}{@{}p{1.8cm}p{2.6cm}p{3.0cm}p{3.6cm}@{}}
    \toprule
    \textbf{Problem} & \textbf{Example} & \textbf{First question} & \textbf{Possible action} \\
    \midrule
    Missing & No ADC reading & Why is it missing? & Leave, remove, or interpolate \\
    Noise & Small random variation & Within expected uncertainty? & Usually retain \\
    Outlier & One extreme reading & Is it physically valid? & Retain/flag/remove --- after investigating \\
    Sensor failure & Impossible value & Is the system faulty? & Remove or mark invalid \\
    Genuine event & Sudden but real change & Physically meaningful? & \textbf{Keep} \\
    Long gap & Sensor offline & Can interpolation be trusted? & Usually do not interpolate blindly \\
    \bottomrule
    \end{tabular}
    \end{center}
\end{frame}

\begin{frame}{Chapter Summary}
    \leavevmode
    \small
    \begin{itemize}
        \item Errors in \textbf{code} are detected by Python; flaws in
              \textbf{data} are not, because a corrupted reading is
              still a valid number.
        \item Catching a \alert{specific} exception beats a bare
              \texttt{except:}, which hides faults you did not
              anticipate.
        \item Never compare floats with \texttt{==} --- use
              \texttt{math.isclose()}.
        \item \texttt{NaN} is contagious: one bad value poisons a mean.
              \texttt{np.isnan()} locates it; NaN-aware functions skip
              it --- skipping is not the same as repairing.
        \item Deciding \alert{what} to do about a gap matters more than
              knowing the functions.
        \item An outlier is a statistical description, not a verdict:
              \alert{Detect $\to$ Flag $\to$ Investigate $\to$ Decide
              $\to$ Document}.
    \end{itemize}

    \vspace{3mm}
    \begin{center}
        \large\color{primary}\textbf{Next: Interpolation, Curve Fitting
        and Calibration}
    \end{center}
\end{frame}

\end{document}
